\documentclass[12pt,a4paper]{article}
\usepackage[margin=1in]{geometry}
\usepackage{graphicx}
\usepackage{float}
\usepackage{array}
\usepackage{longtable}
\usepackage{xcolor}
\usepackage{hyperref}
\usepackage{listings}
\usepackage{titlesec}

\hypersetup{
    colorlinks=true,
    linkcolor=blue,
    urlcolor=blue
}

\titleformat{\section}{\large\bfseries}{\thesection.}{0.5em}{}
\titleformat{\subsection}{\normalsize\bfseries}{\thesubsection.}{0.5em}{}

\lstset{
    basicstyle=\ttfamily\small,
    breaklines=true,
    frame=single,
    backgroundcolor=\color{gray!8},
    columns=fullflexible
}

\title{\textbf{Bean Management in Spring}\\[4pt]\large Assignment Report}
\author{Student Submission}
\date{\today}

\begin{document}

\maketitle

\section{Overview}
This report presents the implementation of the Maven project \texttt{chapter02\_assignment} for the topic \textbf{Management of Beans in Spring}. The solution uses annotation-based configuration with Spring components, service and repository layers, and bean lifecycle management.

\section{Project Objective}
The assignment required the following:
\begin{itemize}
    \item Create a Maven project named \texttt{chapter02\_assignment}.
    \item Implement \texttt{StudentDao} using \texttt{@Repository}.
    \item Implement \texttt{StudentService} using \texttt{@Service} and inject \texttt{StudentDao} with \texttt{@Autowired}.
    \item Implement \texttt{CourseInfo} using \texttt{@Component}, lifecycle annotations, and bean scope testing.
    \item Create \texttt{AppConfig} using Java configuration and component scanning.
    \item Create \texttt{SpringApp} to initialize the Spring container and test bean behavior.
\end{itemize}

\section{Project Structure}
\begin{lstlisting}
chapter02_assignment
|-- pom.xml
|-- src
    |-- main
        |-- java
            |-- com.example.chapter02
                |-- SpringApp.java
                |-- component
                |   |-- CourseInfo.java
                |-- config
                |   |-- AppConfig.java
                |-- dao
                |   |-- StudentDao.java
                |-- service
                    |-- StudentService.java
\end{lstlisting}

\section{Implementation Details}

\subsection{StudentDao}
The \texttt{StudentDao} class is annotated with \texttt{@Repository}. It provides the method \texttt{saveStudent()} which prints a success message when student data is saved.

\begin{lstlisting}[language=Java]
@Repository
public class StudentDao {
    public void saveStudent() {
        System.out.println("StudentDao: student information saved successfully.");
    }
}
\end{lstlisting}

\subsection{StudentService}
The \texttt{StudentService} class is annotated with \texttt{@Service}. It uses \texttt{@Autowired} to inject \texttt{StudentDao}. The method \texttt{registerStudent()} prints a registration message and then calls \texttt{studentDao.saveStudent()}.

\begin{lstlisting}[language=Java]
@Service
public class StudentService {

    @Autowired
    private StudentDao studentDao;

    public void registerStudent() {
        System.out.println("StudentService: registering student...");
        studentDao.saveStudent();
    }
}
\end{lstlisting}

\subsection{CourseInfo}
The \texttt{CourseInfo} bean is annotated with \texttt{@Component}. It demonstrates:
\begin{itemize}
    \item constructor execution,
    \item initialization with \texttt{@PostConstruct},
    \item destruction with \texttt{@PreDestroy},
    \item bean scope testing using \texttt{@Scope("singleton")} and later \texttt{@Scope("prototype")}.
\end{itemize}

\begin{lstlisting}[language=Java]
@Component
@Scope("singleton")
public class CourseInfo {

    private String courseName = "Spring Framework Basics";

    public CourseInfo() {
        System.out.println("CourseInfo constructor called.");
    }

    @PostConstruct
    public void init() {
        System.out.println("CourseInfo initialized.");
    }

    @PreDestroy
    public void destroy() {
        System.out.println("CourseInfo destroyed.");
    }

    public String getCourseName() {
        return courseName;
    }
}
\end{lstlisting}

\subsection{AppConfig}
The configuration class uses annotation-based component scanning:

\begin{lstlisting}[language=Java]
@Configuration
@ComponentScan("com.example.chapter02")
public class AppConfig {
}
\end{lstlisting}

\subsection{SpringApp}
The main class initializes the Spring container, retrieves beans, checks the scope behavior, prints the course name, and finally closes the container.

\begin{lstlisting}[language=Java]
public class SpringApp {
    public static void main(String[] args) {
        AnnotationConfigApplicationContext context =
                new AnnotationConfigApplicationContext(AppConfig.class);

        StudentService studentService = context.getBean(StudentService.class);
        studentService.registerStudent();

        CourseInfo course1 = context.getBean(CourseInfo.class);
        CourseInfo course2 = context.getBean(CourseInfo.class);

        System.out.println("course1 == course2 : " + (course1 == course2));
        System.out.println(course1.getCourseName());

        context.close();
    }
}
\end{lstlisting}

\section{Observed Output}
The application was executed successfully with \texttt{@Scope("singleton")}. The terminal output was:

\begin{lstlisting}
CourseInfo constructor called.
CourseInfo initialized.
StudentService: registering student...
StudentDao: student information saved successfully.
course1 == course2 : true
Spring Framework Basics
CourseInfo destroyed.
\end{lstlisting}

\section{Scope Test Analysis}
\subsection{Singleton Scope}
When \texttt{@Scope("singleton")} is used:
\begin{itemize}
    \item only one object of \texttt{CourseInfo} is created,
    \item both \texttt{course1} and \texttt{course2} refer to the same object,
    \item therefore the output is:
\end{itemize}

\begin{lstlisting}
course1 == course2 : true
\end{lstlisting}

\subsection{Prototype Scope}
When \texttt{@Scope("prototype")} is used:
\begin{itemize}
    \item a new object is created every time the bean is requested,
    \item \texttt{course1} and \texttt{course2} become different objects,
    \item therefore the output changes to:
\end{itemize}

\begin{lstlisting}
course1 == course2 : false
\end{lstlisting}

\section{Screenshots}
Figure~\ref{fig:output} is intended for the terminal output screenshot provided with the assignment evidence.

\begin{figure}[H]
    \centering
    \includegraphics[width=\textwidth]{output-screenshot.png}
    \caption{Terminal output of the Spring application execution}
    \label{fig:output}
\end{figure}

\noindent\textbf{Note:} Save the screenshot you shared as \texttt{output-screenshot.png} in the same folder as \texttt{percent.tex} before compiling the report.

\section{Conclusion}
The assignment was completed successfully using annotation-based bean configuration in Spring. The implementation demonstrates repository and service bean management, dependency injection with \texttt{@Autowired}, bean lifecycle methods with \texttt{@PostConstruct} and \texttt{@PreDestroy}, and scope behavior testing with singleton and prototype beans.

\end{document}
